\subsection*{Wozu eine Abschlussarbeit?}

Mit der Abschlussarbeit soll der oder die Studierende nachweisen, dass er oder sie in begrenzter Zeit mit den einschlägigen Methoden der betreffenden Fachrichtung komplexere Themen selbständig bearbeiten und Aufgaben lösen kann \cite{RSTPO_HTW}. Ein wesentlicher Bestandteil der Abschlussarbeit ist die schriftliche  wissenschaftiche Darstellung der im Rahmen der Abschlussarbeit durchgeführten Arbeiten. 

\begin{itemize}
    \item Auf der Internetseite \hyperlink{
https://www.f1.htw-berlin.de/studieren/abschlussarbeit/}{\emph{Abschlussarbeiten am FB~1}} befinden sich Antworten zum  formalen Ablauf einer Abschlussrabeit. Fragen wie ''Wie beantrage ich eine Zulassung zur Abschlussarbeit?'' oder ''Kann man die Bearbeitungszeit verlängern?'' werden hier bantwortet.
    \item Viele Tipps zum Organisieren der Abschlussarbeit findest du im Moodlekurs\\ 
    \hyperlink{moodle.htw-berlin.de/course/view.php?id=17639}{\emph{Infoportal wissenschaftliches Schreiben}}.
    \item Der offizielle Leitfaden vom FB 1 ist\\\hyperlink{https://www.f1.htw-berlin.de/fileadmin/HTW/Zentral/FB/FB1/Leitfaden_zur_Erstellung_einer_Wiss.-Arbeit_neu.pdf}{\emph{Leitfaden zur Erstellung einer wissenschaftlichen Arbeit am Fachbereich 1}}.\cite{fb1leitfaden}.
    \item In \cite{volkerformeln} wird ausführlich erlärt, wie du Formelzeichen richtig setzt.
\end{itemize}

Dieses Dokument nimmt Bezug auf die genannten Dokumente. Es bündelt alles Wesentliche in einem Dokument möglichst in Form von Checklisten. Außerdem dient dieses Dokument als Vorlage für die schriftliche Ausarbeitung.
Die Kapitelüberschriften sollten für die Abschlussarbeit beibehalten werden.

Zunächst werden in  Kapitel~\ref{sec:methoden}  wichtige Punkte zum wisschenschaftlichen Schreiben herausgestellt. Es geht um Formulierungen in Wissenschftlöivrn Arbeiten im allgemeinen \ref{sec:writing} und die Darstellung mit Bilfern und Tabellen (\ref{sec:bildertabs} Insbesondere wird die korrekte Darstellung mathematrischer Sachverhalte erläutert \ref{sec:formulas}, Abschnitt \ref{sec:zitieren} reklart das richtige Zitieren und Abschnitt~\ref{sec:lizenzen} wie urheberrechtlichr Schutz gewärleistet weden kann.
%\subsection{schriftliche Ausarbeitung}

\subsection{Umfang}

Der Abschlussbericht zur Arbeit ist ein technischer Bericht. Er soll eine geordnete Darstellung der im Verlauf der Arbeit angestellten Überlegungen, durchgeführten Ableitungen, entwickelten Algorithmen, Experimenten usw. im Sinne einer Diskussion bringen. Das Herausarbeiten und Bewerten der Ergebnisse ist dabei die we-sentliche geistige Leistung. Eine konzentrierte Darstellung (ca. 30 bis 50 A4-Seiten, 1 1/2-zeilig = 4 Zei-len/Zoll) ist anzustreben; vermeiden Sie unbedingt weitschweifige Abhandlungen. Experimentprotokolle, Rechnerausdrucke, etc. können geordnet, eindeutig und vollständig beschriftet als Anhang beigeheftet wer-den. Berechnungen sind trotz der gebotenen Kürze so ausführlich zu bringen, sodass eine schrittweise Überprüfung möglich ist.

\subsection{Äußere Form}

Die Arbeit ist in Deutsch oder Englisch abzufassen und mit dem Textverarbeitungssystem TeX zu schreiben. Alle Berichte sind doppelseitig und in schwarz/weiß auszudrucken. Studien- und Bachelorarbeiten werden mit zwei Tackern an der linken Seite gebunden. Bei Diplom- oder Masterarbeiten ist eine Version gebunden (Klebebindung, keine Ringbindung!) abzugeben, Zwischenberichte reichen mit Tackern. 
Am Leibniz-Rechenzentrum (LRZ) sind PCs und Drucker zum Erstellen der Diplomarbeit verfügbar. Eine LRZ-Benutzerkarte erhalten Sie bei Frau Renner-Šmid (Zi. N2515).


\subsection{Datenablage}
Zur Archivierung der Arbeit ist die Erstellung einer CD oder DVD erwünscht, die 

1)	alle Quellen, die zur Erstellung der Berichte benutzt worden sind (Das sind also zum Beispiel: .tex-File, .bib-File, Stylefiles, alle Graphiken, bitte mit absoluten/ relativen Pfaden aufpassen, dass alles ohne Fehler kompilierbar ist...)
2)	die gehaltenen Präsentationen, wieder mit allen benötigten Quellen wie Graphiken, .tex-Files, ppt-File...
3)	alle Daten, die zum Reproduzieren der Ergebnisse benötigt werden. (z.B. Simulationen, die im Be-richt gezeigten Daten, Daten zu Messkurven)
enthält.

Eine CD/DVD muss nur einem Exemplar beigelegt werden.

Zeichnungen sind unter Verwendung genormter Schaltzeichen, Bildzeichen und Symbole anzufertigen und zu beschriften. Dies gilt besonders auch für Diagramme. Zeichenblätter, deren Format größer als DIN A4 ist, sind gefaltet dem Text beizuheften oder beizulegen.


\subsection{beispil struktur}

\begin{itemize}
    \item Titelseite gemäß Vorlage.
    \item Kurzfassung (Abstract) in Deutsch und in Englisch (max. je 15 Zeilen,  vgl. Muster).
    \item Aufgabenbeschreibung im Original.
    \item Inhaltsverzeichnis der Diplomarbeit mit Seitenangaben.
    \item Vollständiges Verzeichnis aller verwendeten Buchstaben und Symbole.
    \item Einleitung (u. a. Einordnung und Detaillierung der Aufgabenstellung).
    \item Hauptteil(e) der Diplomarbeit.
    \end{itemize}
    \item Zusammenfassung (ca. 1 - 2 Seiten): Zusammenstellung und kritische Bewertung der wesentlichen Ergebnisse., Ausblick über Möglichkeiten zur Verbesserung und Weiterarbeit.
    \item Verzeichnis der bei der Anfertigung der Diplomarbeit verwendeten Literatur. Die Literaturhinwei-se sind im Text an den entsprechenden Stellen anzugeben.
    \item Anhang: Nicht im Text eingefügte Bilder und Tabellen sind am Schluss der Arbeit mit fortlaufender Numme-rierung anzufügen., Programmauflistungen und -dokumentationen (diese können ggf. auch in einem zweiten Band zu-sammengestellt werden).
\end{itemize}
	
Eine Tex-Vorlage in OVERLEAF!


\section{schriftlihce ausarbeitung Struktur}
Die Dokumentation sollte zunächst dem gebräuchlichen Standard für wissenschaftliche Arbeiten folgen.
Hinweise finden sich auf den Web-Seiten der HTW-Berlin. Nachfolgend wird auf einzelne Aspekte eingegangen,
die beim Anfertigen der Dokumentation beachtet werden sollten.

Äußere Form der Dokumentation
Deckblatt. Auf dem Deckblatt sollte der angemeldete Titel der Abschlussarbeit, Ihr vollständiger Name,
Ihre Matrikelnummer, das Datum der Abgabe und die Namen und Titel der Gutachter angegeben werden
(siehe Anhang). Üblich ist außerdem die Verwendung des HTW-Logos. Falls die Arbeit mit einem
Sperrvermerk versehen ist, muss auf dem Deckblatt gut sichtbar der Hinweis "Außerhalb des Prüfungsverfahrens unzugänglich" angebracht sein. Alternativ können Sie den Vermerk auch auf einer separaten
Seite platzieren, die dann unmittelbar vor oder nach dem Deckblatt einzubinden ist.
Eidesstattliche Erklärung. Sie müssen Ihrer Dokumentation eine unterzeichnete eidesstattliche Erklärung
beifügen. Darin müssen Sie erklären, dass Sie selbständig und ohne fremde Hilfe gearbeitet, alle von
Ihnen genutzten Quellen und Hilfsmittel angegeben, Zitate gekennzeichnet und die Abschlussarbeit bis her keiner anderen Prüfbehörde vorgelegt haben. Ein Muster für eine eidesstattliche Erklärung finden Sie
im Anhang oder unter [7].

Es ist üblich, die Nummerierung der Absätze auf drei Ebenen zu beschränken (also z.B. Kapitel 1 mit Unterkapitel
1.2 und „Unter-Unterkapitel“ 1.2.3).

Weiterhin sollten Sie bei allen Verzeichnissen auf Vollständigkeit achten. Dies betrifft vor allem das Abkürzungsverzeichnis,
da es des Öfteren vorkommt, dass verwendete Abkürzungen vom Autor übersehen
werden. Solche Lücken fallen vor allem bei weniger gebräuchlichen Abkürzungen schnell auf.

Quellenverzeichnis (Literaturverzeichnis, Webverzeichnis). Eine wissenschaftliche Arbeit ist immer eingebettet
in einen Kontext existierenden Wissens, der unbedingt zu dokumentieren ist. Jede Behauptung,
jedes nichtoffensichtliche Detail, jedes Zitat, jedes nicht selbst angefertigte Bild usw. muss mit einer
Quellenangabe versehen sein. Wenn Sie z.B. schreiben wollen, dass in Vatikanstadt 2,27 Päpste pro
Quadratkilometer leben, sollten Sie eine Quellenangabe beifügen, die auf die entsprechende Statistik
verweist (als Beleg der Behauptung). Falls Sie selber herausgefunden haben, dass in Vatikanstadt 2,27
Päpste pro Quadratkilometer leben, müssen Sie als Quelle zumindest angeben, woher Sie die Information
haben, dass Vatikanstadt 0,44 km² groß ist (als Beleg eines Fakts) [8].

\section{schriftliche ausarbeitung inhalt}


\subsection{Formeln}#



\sebsection{Bilder& Abbildungen}
}Bilder/Abbildungen. Bilder sind in einer Ingenieursarbeit von essentieller Bedeutung. Jedem wird klar
sein, dass es kaum Sinn macht, einen Schaltplan, eine Konstruktionszeichnung, einen Ablaufplan usw.
ausschließlich verbal zu beschreiben. Dabei müssen Bilder keinem vorgegebenen Standard folgen, sondern
lassen sich, wie die Sprache, als ein freies Ausdrucksmittel nutzen. Im Vergleich zur Sprache haben
Bilder aber den Vorteil, dass Sie wenigstens zwei Dimensionen besitzen. Durch eine perspektivische Darstellung
und durch Verwendung unterschiedlicher Farben und Schattierungen lassen sich weitere Dimensionen
ergänzen.
Das Einfügen von vorhandenen Bildern wird in den meisten Textverarbeitungsprogrammen gut unterstützt.
Deshalb machen Bilder aus fremden Quellen kaum Arbeit. Deutlich aufwändiger ist es, eigene
Bilder zu erstellen. Je nach verwendeter Textverarbeitung kann es erforderlich sein, zunächst ein Graphikprogramm
zu starten, das Bild zu erstellen (was Fertigkeiten verlangt), das Bild zu speichern und
anschließend in die Dokumentation einzubinden. Wegen dieses Aufwands kommt es immer wieder vor,
dass Autoren auf sinnvolle Bilder verzichten und versuchen, einen Sachverhalt verbal zu beschreiben.
Das Ergebnis ist normalerweise unbefriedigend und wird daher auch entsprechend bewertet.


\subsection{Tabellen}
Tabellen. Tabellen sind mit Bildern vergleichbar, weshalb das zuvor Gesagte weitgehend auch für Tabellen
gilt. Es gibt wenige Abweichungen, auf die hier kurz hingewiesen werden soll. Normalerweise verwendet
man bei Tabellen nicht eine Tabellenunterschrift, sondern eine der Tabelle vorangestellte Überschrift.
Genau wie Bilder und Grafiken müssen auch Tabellen nummeriert werden. Üblich ist es, für Tabellen
eine eigene Nummerierung zu verwenden, d.h. in einer Dokumentation gibt es möglicherweise
eine Abbildung 1 und eine Tabelle 1.

\subsection{Text}

Text. Bei der Formulierung eines Textes gilt zunächst genau das, was bereits in der Schule gelernt wurde.
Vermeiden Sie Rechtschreib- und Zeichensetzungsfehler. Bei der Erkennung von Rechtschreibfehler werden
Sie durch moderne Textverarbeitungssoftware hervorragend unterstützt. Zeichensetzungsfehler
werden jedoch meist nicht automatisch erkannt. Hier ist es wichtig, dass Sie Ihre Arbeit einmal selbst
korrekturlesen und im Idealfall ein zweites Mal durch jemanden, der mit der deutschen Zeichensetzung
gut vertraut ist. Ihnen wird verziehen, wenn Sie nicht in der Lage sein sollten, eine Differenzialgleichung
zu lösen (was ein Ingenieur vielleicht beherrschen sollte). Ihnen verzeiht jedoch niemand, wenn Sie einen
Text schreiben, der voller Rechtschreibfehler ist (was nicht die Kernkompetenz eines Ingenieurs ist).
Des Weiteren sollten Sie auf grammatikalische Fehler achten: So sind z.B. Sätze ohne Verb im Deutschen
nicht möglich, treten aber immer wieder in Erscheinung. Vermeiden Sie Wortwiederholungen. Berücksichtigen
Sie jedoch: In wissenschaftlich technischen Bereich sind Fachbegriffe oft eng umgrenzt. Daher
kann die Verwendung von unterschiedlichen Fachbegriffen für einen speziellen Sachverhalt zu Missverständnissen
führen.
Wenn Sie z.B. in Ihrer Dokumentation einmal von einer Speicherverwaltungseinheit und einmal von einer
MMU reden, weiß der Leser nicht, ob Sie dasselbe oder zwei unterschiedliche Dinge meinen. Hierbei hilft
es nicht, einleitend zu erläutern, dass die beiden Begriffe synonym zu verstehen sind. Falls der Leser diese
kurze Begriffserklärung nämlich überliest, weil er für einen Moment unaufmerksam ist, ärgert er sich
anschließend mit Ihrem Text herum. Übrigens lassen sich Wortwiederholungen bei Fachbegriffen immer
durch geschicktes Formulieren vermeiden.

Genauso wie synonyme Fachbegriffe sollten Sie ungebräuchliche Abkürzungen vermeiden, und zwar
auch dann, wenn Sie diese in Ihrer Dokumentation explizit in einem Satz erklären. Es ist nicht ganz eindeutig,
wann eine Abkürzung gebräuchlich oder nicht gebräuchlich ist. Unproblematisch ist die Verwendung
von Abkürzungen, die jedermann versteht, wie z.B. PKW, FAX, USB und PC. Die Abkürzungen ET,
FFT und CPU sind hingegen weniger verbreitet, werden aber von Elektrotechnikern verstanden. Sie dürfen
in einer wissenschaftlich technischen Dokumentation des Studiengangs Elektrotechnik deshalb eingeschränkt
zum Einsatz kommen. Verzichten Sie aber auf Abkürzungen wie CRT, SSE oder TGS1. Sie werden
möglicherweise in der Firma verstanden, in der Sie Ihre Abschlussarbeit anfertigen, aber vielleicht
nicht mehr von Ihren Gutachtern.

\subsection{Struktuir}

\subsection{rest}

\begin{table}[ht]
	%\centering\renewcommand{\arraystretch}{1.1} 
	\footnotesize
	\begin{tabularx}{0.9\textwidth}{llXXXXrr}
		\midrule
		\multicolumn{2}{c}{Bruchpunkte} & \multicolumn{4}{c}{Strukturbruchtermine} & 					\multicolumn{1}{c}{BIC}& \\
		\midrule
		\multicolumn{2}{c}{0}   				&       	&   	    &       	&	      	&	 -2.337& \\
		\multicolumn{2}{c}{1}    				& 03/2009 &     	  & 	      & 	      &  -2.369& \\
		\multicolumn{2}{c}{~\,2$^\star$}& 04/2007 & 03/2009 &   &   	    &  -2.395& \\
		\multicolumn{2}{c}{3}    				& 04/2006 & 03/2009 & 11/2010 &     	  &  -2.373& \\
		\multicolumn{2}{c}{4}    				& 07/2006 & 01/2008 & 07/2009 & 01/2011 &  -2.313& \\
		\midrule\\[-3mm]
		\multicolumn{8}{l}{\begin{minipage}{0.88\textwidth}Ergebnisse des CUSUM-Tests für monatliche \textsl{4TC+2Cal} FFA Log-Differenzen von 01/2005 - 12/2012. Die BIC-optimale Spezifikation ist durch $^\star$ gekennzeichnet.\end{minipage}}
	\end{tabularx}
	\caption{Ergebnisse des CUSUM Strukturbruchtests}
	\label{tab:cusum}
\end{table}



Ergebnis der Linearisierung eines nichtlinearen Systems um eine zeitunabhängige Referenzlösung  ist
ein LTI-Zustandsmodell in der Standardform 
%
\begin{equation}
\begin{array}{rcl}
\underline{\dot{x}} & = & A\underline{x} + B\underline{u}\\
\underline{y} & = & C\underline{x} + D\underline{u}
\end{array}\,.
\label{eq:linear}
\end{equation}
%
mit auf die
Ruhelage bezogenen konstanten Zu\-stands\-matrizen.
\begin{equation}
\begin{array}{rcl}
\Delta\ul{\dot{x}} & = & A\Delta\ul{x}+B\Delta\ul{u}\\
\Delta\ul{y} & = & C\Delta\ul{x}+D\Delta\ul{u}
\end{array}
\label{eq:1.11nochmal}
\end{equation}

Hinweis: Geht aus dem Zusammenhang hervor, dass es sich bei \mref{1.11nochmal} um ein
Klein\-sig\-nal\-mo\-dell handelt, dann wird zur Vereinfachung der
Schreibweise das Symbol $\Delta$ häufig unterdrückt.


\subsection{Bilder}
\label{sec:bilder}


\begin{figure}[!ht]
	\centering
\psfrag{#Wf}[lc][rb]{\footnotesize \emph{(reference value)}}
\psfrag{#Wd}[lc][lc]{\small  Führungsgröße $w(t)$}
\psfrag{#Zf}[lc][lc]{\footnotesize \emph{(disturbance)}}
\psfrag{#Zd1}[lc][lc]{\small Störgröße $z_1(t)$}
\psfrag{#Zd2}[lc][lc]{\small Störgröße $z_2(t)$}
\psfrag{#Zd3}[lc][lc]{\small Störgröße $z_3(t)$}
\psfrag{#Ee}[lc][lc]{\footnotesize \emph{(closed loop error)}}
\psfrag{#Ed}[lc][lc]{\small Regelabweichung $e(t)$}
\psfrag{#Ue}[lc][lc]{\footnotesize \emph{(control output)}}
\psfrag{#Ud}[lc][lc]{\small Reglerausgang $u(t)$}
\psfrag{#Ua}[lc][lc]{\small Stellgröße $u_a(t)$}
\psfrag{#Uae}[lc][lc]{\footnotesize \emph{(manipulated variable)}}
\psfrag{#Y2e}[lc][lc]{\footnotesize \emph{(control variable)}}
\psfrag{#Y2d}[lc][lc]{\small Regelgröße $y(t)$}
\psfrag{#Y1e}[lc][lc]{\footnotesize \emph{(feedback variable)}}
\psfrag{#Y1d}[lc][lc]{\small Rückführgröße $y_s(t)$}
\psfrag{#Str}[cc][cc]{\BildFont {\large  	 Regelstrecke}}
\psfrag{#R}[cc][cc]{\BildFont Regler}
\psfrag{#Sen}[cc][cc]{\BildFont Sensor}
\psfrag{#P}[cc][cc]{\BildFont System}
\psfrag{#A}[cc][cc]{\BildFont Aktor}
\hspace*{-2cm}
	\includegraphics[width=0.9\textwidth]{regelkreis3.eps}
	\caption{Einschleifiger Regelkreis.}
	\label{fig:regelkreis}
\end{figure}

%Formelschreibweise und Formelsatz
in wissenschaftlichen Arbeiten
Gültige Normen
Zur Schreibweise von Formeln gibt es gültige Normen, die in wissenschaftlichen
Arbeiten an der HTW Berlin zu beachten sind:
 DIN 1338 Formelschreibweise und Formelsatz
 DIN 1304 Formelzeichen
 DIN 1301 Einheitennamen, Einheitenzeichen
 DIN 1313 Größen
Formelsatz
Die Normen regeln klar, welche Buchstaben für Größen und Einheiten zu verwenden
und wie Formeln darzustellen sind. Dabei ist auch geregelt, welche
Zeichen gerade und welche kursiv zu setzen sowie wo Leerzeichen einzufügen
sind.
Es existiert der allgemeine Glaube, dass Formeleditoren von Textverarbeitungsprogrammen
wie MS-Word automatisch die Formeln korrekt setzen. Wer sich
darauf verlässt, wird einiges erreichen. Korrekt geschriebene Formeln werden
aber sicher nicht dazu gehören. Der Formelsatz muss stets manuell angepasst
werden.
Die wichtigsten Regeln für einen korrekten Formelsatz sind im Folgenden zusammengefasst.
 Variablen und Größen werden kursiv gesetzt
 Zahlen und Einheiten werden gerade gesetzt
 Zwischen Variable und Einheit steht immer ein Leerzeichen. Ausnahmen: Beim
Prozent-Zeichen (%) kann das Leerzeichen wahlweise stehen oder entfallen.
Bei Winkelbezeichnungen und Gradangaben (°) entfällt das Leerzeichen.
Beispiele
�� �� 220 V
�� �� �� ∙ �� �� 220 V ∙ 5 A �� 1100 W �� 1,1 kW
�� �� 5 %
�� �� 5,7°
 Indizes werden gerade gesetzt, wenn sie die Variable näher bezeichnen. Lediglich
variable Indizes werden kursiv gesetzt.
Beispiele für bezeichnende Indizes
���� �� 37,5 V (Leerlaufspannung)
�������� �� 99,9 W (MPP-Leistung)
���� ������ ������������ �� 14° (maximale Sonnenhöhe im Winter)



Formelschreibweise und Formelsatz in wissenschaftlichen Arbeiten 2
Beispiel für bezeichnende und variable Indizes
�������� �� ���� ��
∙ ����
��
������
 Die mathematischen Konstanten , e und i, mathematische Funktionen wie
sin, cos, tan, lim, ln, exp sowie Summen- und Integralzeichen werden gerade
gesetzt.
Beispiele
�� �� π ����
���� ��
1
sin ����
���� �� ���� ∙
ln ����
����
ln ����
����
 Chemische Formelzeichen und Abkürzungen werden gerade gesetzt.
Beispiele
H��SO�� (Schwefelsäure)
C��H��O (Ethanol)
WKA (Windkraftanlage)
MPP (Maximum Power Point)
Einheitenangaben bei Tabellen und Grafiken
Bei der Verwendung von eckigen Klammern im Zusammenhang mit Einheiten
existiert in zahlreichen Dokumenten, Skripten und gar Fachbüchern ein regelrechtes
formales Chaos. Dabei ist die Verwendung von eckigen und geschweiften
Klammern in den entsprechenden Normen und der SI-Einheitendefinition eindeutig
geregelt. Diese Regeln sollen auch bei wissenschaftlichen Arbeiten an der
HTW berücksichtigt werden. Danach gilt:
Größe = {Größe} · [Größe]
Dabei bedeuten die eckigen Klammern „Einheit von“ und die geschweiften Klammern
„Zahlenwert von“.
Beispiele
�� �� 590 W
������ �� W
������ �� 590
Bei Beachtung der Regeln ergeben sich drei formal korrekte Beschriftungen von
Koordinatensystemen in Diagrammen.

Beispiele für korrekte Achsenbeschriftungen
Beispiel für formal falsche Achsenbeschriftung
Prof. Dr. Volker Quaschning
2012-11-11
0
1
2
3
4
5
6
7
8
9
0 5 10 15 20 25 30 35 40 45
Modulstrom I in A
Modulspannung U in V
0
1
2
3
4
5
6
7
8
9
0 5 10 15 20 25 30 35 40 45
Modulstrom I/A
Modulspannung U/V
0
1
2
3
4
5
6
7
8
9
0 5 10 15 20 25 30 35 40 45
Modulstrom I
Modulspannung U
A
V
0
1
2
3
4
5
6
7
8
9
0 5 10 15 20 25 30 35 40 45
Modulstrom I [A]
Modulspannung U [V]
%
\subsection{Zitieren}
\label{sec:zitieren}

\begin{itemize}
\item Wenn Sie ein Dokument verfassen, das in irgendeiner Weise verbreitet wird (oder werden könnte), sollten Sie darauf achten, dass alle Ihre Abbildungen frei von Urheberrechtsproblemen sind. Das bedeutet: Entweder Sie zeichnen sie selbst (auch die Übernahme von Teilen aus anderen Bildern und deren Veränderung ist ein Problem) oder Sie holen eine Genehmigung ein und geben in der Bildunterschrift an, dass Sie diese Genehmigung erhalten haben und von wem. Meistens sagen Ihnen die Verlage, wie man das schreibt, und haben sogar ein automatisiertes Verfahren, um die Erlaubnis zu erhalten, z. B. über RightsLink/Copyright Clearance Center (meist sofort und kostenlos für Abschlussarbeiten von Studierenden). Auch Bilder aus dem öffentlichen Bereich müssen ordnungsgemäß referenziert werden. Es gibt auch verschiedene Arten von Lizenzen, z. B. erlauben einige keine Änderungen an einem Bild.

\item Vermeiden Sie direkte Zitate (``...'') wann immer möglich, formulieren Sie stattdessen Aussagen um, wenn Sie zitieren.

\item Es ist eindeutig vorzuziehen, Primär- und Originalliteratur zu zitieren, nicht Rezensionen oder Lehrbücher. Gehen Sie auf die Originalquellen zurück. Wenn Sie Rezensionen oder Bücher zitieren, tun Sie dies ausdrücklich, z. B. mit ``ein Überblick findet sich in...''.

\item  \item Lesen Sie jede Arbeit, die Sie zitieren, so lange, bis Sie sie verstehen, oder zumindest den Teil der Arbeit, auf den Sie sich beziehen. Verlassen Sie sich nicht nur auf Zusammenfassungen oder, schlimmer noch, auf Sekundärliteratur, die die Originalquelle (oft falsch) zitiert. Wenn Sie es für unmöglich halten, die zitierten Quellen zu lesen und zu verstehen, schränken Sie den Umfang Ihrer Arbeit so weit ein, bis es möglich ist.

\item Belegen Sie jede einzelne Behauptung, die keine offensichtliche Tatsache ist, entweder durch einen Verweis oder durch ein eigenes Experiment, eine Schlussfolgerung oder eine Berechnung. 

\cite{Menge}
\end{itemize}
%\subsection{Rechtevergabe}
\label{sec:lizenzen}

Die Werke eines Schöpfers (wie Texte, Musikstücke, Bilder, Videos usw.) sind normalerweise urheberrechtlich geschützt. Der Schöpfer kann aber entscheiden, dass er Werke anderen Menschen zur Verfügung stellt, ohne dass sie ausdrücklich um Erlaubnis fragen müssen. Dazu veröffentlicht er die Werke mit einem entsprechenden Hinweis, dass er zum Beispiel das Recht zum Kopieren, Verändern und Wiederveröffentlichen allen anderen zugesteht.

Für juristische Laien ist es allerdings schwierig, einen entsprechenden Rechtstext zu formulieren. Schließlich soll deutlich sein, was erlaubt ist und was nicht, und es soll auch kein Missbrauch mit den zur Verfügung gestellten Werken möglich sein (etwa, dass jemand behauptet, er selbst sei Schöpfer dieser Werke). Um diesem Problem zu begegnen, wurde die Organisation Creative Commons gegründet, um solche Rechtstexte (Lizenzen) zu erarbeiten. \cite{wikicc}



\hyperlink{https://www.wur.nl/en/article/What-are-Creative-Commons-licenses.htm}{FAQ- What are Creative Commons licenses?}

\hyperlink{https://creativecommons.org/licenses/}{CC Bilder }

\hyperlink{https://moqod-software.medium.com/understanding-open-source-and-free-software-licensing-c0fa600106c9}{Software tetxte}

\hyperlink{https://choosealicense.com/}{guide}

In Kapitel~\ref{sec:ergebnisse} stelle ich die drei gedforderten Ergebnisse der Abschlusarbeit vor, die schriftliiche Ausarbeitung, das Kolloquium und die Daten.
%In \ref{sec:kolloquium}
%\ref{sec:latexinstallation}
%\ref{sec:writing}
In Kapitel~\ref{sec:diskussion} zeigt die Bewertungskriterien

Im Anhang befinden sich noch viele Ergänzungen. 















